\documentclass[11pt,a4paper]{article}

\usepackage[utf8]{inputenc}
\usepackage[T1]{fontenc}
\usepackage{lmodern}
\usepackage{amsmath,amssymb,amsthm}
\usepackage{booktabs}
\usepackage{graphicx}
\usepackage{hyperref}
\usepackage{xcolor}
\usepackage{listings}
\usepackage{geometry}
\usepackage{enumitem}
\usepackage{caption}
\usepackage{array}
\usepackage{float}

\geometry{margin=1in}

\hypersetup{
    colorlinks=true,
    linkcolor=blue!70!black,
    citecolor=green!50!black,
    urlcolor=blue!60!black
}

\lstset{
    basicstyle=\ttfamily\small,
    breaklines=true,
    frame=single,
    backgroundcolor=\color{gray!5},
    xleftmargin=1em,
    xrightmargin=1em,
}

\newtheorem{theorem}{Theorem}
\newtheorem{lemma}[theorem]{Lemma}

\newcommand{\R}{\mathbb{R}}
\newcommand{\bO}{\mathcal{O}}
\newcommand{\norm}[1]{\left\|#1\right\|}

\begin{document}

\begin{center}
    {\LARGE\bfseries AXIM: A Universal CUDA-Free Compute and Graphics Runtime}\\[0.5em]
    {\large A Vendor-Neutral Execution Substrate for AI Inference, HPC, and\\
    Graphics, with SYNAXIM as a Case Study}\\[1.5em]
    {\large GRRNMAKER}\\[0.3em]
    {\normalsize \url{https://github.com/GRRN-MAKER/Aximcomp}}\\[0.3em]
    {\normalsize July 2026}
\end{center}

\vspace{1em}\hrule\vspace{1.5em}

\begin{abstract}
The dominant substrate for GPU computing---NVIDIA CUDA---creates a single-vendor dependency
that fragments the accelerator ecosystem. Existing alternatives translate or intercept CUDA
(AMD ROCm/HIP, Intel SYCL, ZLUDA) and therefore remain tethered to CUDA's programming model
and, in most cases, to a single hardware vendor. We present \textbf{AXIM}, a compute and
graphics runtime that \emph{never depends on CUDA at all}. AXIM defines a compact
intermediate representation (the AXIM IR) over a fixed operator set and lowers it to
vendor-neutral targets: AVX-512/AVX2/NEON SIMD on CPUs and Vulkan (SPIR-V) / Metal (MSL) on
GPUs, with runtime device selection. A single kernel definition therefore executes across
NVIDIA, AMD, Intel, and Apple silicon for AI inference, HPC, and games. As a motivating case
study we run SYNAXIM, a framework-free LLM inference engine, entirely on AXIM---including its
INT4 Boolean-bit matvec and low-rank associative memory. We report results
\emph{verified on Apple~M3} (macOS, arm64): the CPU (NEON) and GPU (Metal) backends produce
\textbf{bit-for-bit identical} INT4 matvec output, a full SYNAXIM layer forward runs
CUDA-free, and 24 automated checks pass. Cross-vendor execution (NVIDIA/AMD/Intel via Vulkan)
is presented as a \emph{projected} result grounded in the shipped, compiling SPIR-V shaders.
AXIM is implemented in ${\sim}3{,}400$ lines across Python, Rust, C++/Objective-C++, and
shader source.
\end{abstract}

\tableofcontents
\newpage

\section{Introduction}

\subsection{The Vendor Lock-In Problem}

By 2026, practical GPU computing is effectively synonymous with NVIDIA CUDA. Deep-learning
frameworks, HPC libraries, and increasingly game engines assume a CUDA toolchain, a CUDA
runtime, and NVIDIA hardware. This coupling has three consequences: (i) workloads cannot
move to AMD, Intel, or Apple accelerators without a rewrite; (ii) a single vendor controls
the pace and pricing of the AI compute market; and (iii) commodity CPUs---which ship wide
SIMD units---are treated as second-class citizens for numerical work.

We ask the dual of the question that motivated SYNAXIM. SYNAXIM asked \emph{what is the
minimal software substrate to run a language model?} Here we ask: \textbf{what is the
minimal, fully vendor-neutral substrate required to execute a GPU/CPU compute kernel---and
can the same kernel run everywhere without CUDA?}

\subsection{Contributions}

\begin{enumerate}
    \item \textbf{The AXIM IR} (\S\ref{sec:ir}): a small, explicit operator set with typed
    buffers and memory spaces, sufficient to express SYNAXIM inference and general compute,
    that lowers uniformly to CPU-SIMD and GPU targets.
    \item \textbf{Uniform vendor-neutral lowering} (\S\ref{sec:backend}): one IR compiled to
    AVX-512/AVX2/NEON, SPIR-V (Vulkan), and MSL (Metal), with runtime device selection and
    no CUDA in any path.
    \item \textbf{SYNAXIM on AXIM} (\S\ref{sec:casestudy}): a complete post-transformer
    inference engine running on the runtime with \emph{CPU\,$=$\,GPU bit-exact} output.
    \item \textbf{An honest evaluation} (\S\ref{sec:eval}): claims are separated into
    \emph{verified on Apple~M3} and \emph{projected} (Vulkan cross-vendor), the latter
    grounded in compiling, shipped SPIR-V shaders.
\end{enumerate}

\section{Background and Related Work}

\textbf{CUDA translation and interception.} AMD ROCm/HIP provides a CUDA-like API and a
source translator (HIPIFY) targeting AMD GPUs. Intel oneAPI/SYCL migrates CUDA to SYCL via
SYCLomatic across multiple devices but retains a large runtime. ZLUDA intercepts compiled
CUDA binaries and re-maps them to AMD/Intel drivers. All three are defined relative to
CUDA and inherit its programming model.

\textbf{Portable GPU standards.} Vulkan compute (with SPIR-V) and Apple Metal (with MSL)
are vendor-neutral GPU interfaces already shipped in drivers across NVIDIA, AMD, Intel, and
Apple hardware. AXIM builds directly on Vulkan and Metal and adds nothing between the kernel
and the driver.

\textbf{CPU SIMD.} Modern CPUs expose wide vector units---AVX-512/AVX2 on x86-64 and NEON on
ARM. For bandwidth-bound, single-token inference these are competitive when weights are
quantized. AXIM treats the CPU as a first-class compute target with runtime ISA detection.

\textbf{Position.} Unlike prior work, AXIM does not translate, migrate, or intercept CUDA.
It compiles one IR straight to vendor-neutral targets, and additionally offers an optional
CUDA-compatibility shim (AXIM-HIP) so existing CUDA sources can be \emph{ported} (one
\texttt{\#include}) rather than the runtime depending on CUDA.

\section{The AXIM Intermediate Representation}\label{sec:ir}

The AXIM IR is deliberately small. A program is a \texttt{Module} of one or more
\texttt{Kernel}s; each kernel operates over typed \texttt{Buffer}s annotated with a data
type (\texttt{DType}) and a memory space (\texttt{MemSpace}). The operator set
(\texttt{OpCode}) is fixed and chosen to be sufficient for SYNAXIM while remaining general:

\begin{lstlisting}
ADD  MUL  SUB  SILU  MATVEC  INT4_MATVEC  RMSNORM
LOWRANK_UPDATE  LOWRANK_RETRIEVE  SWIGLU
\end{lstlisting}

Two properties make the IR portable. First, \textbf{explicit typing and memory space}: every
buffer declares its \texttt{DType} and \texttt{MemSpace}, so the same IR lowers to a CPU
(host memory, SIMD loop) or a GPU (device buffer, dispatch) without changing operator
semantics. Second, \textbf{backend-agnostic operators}: operators describe \emph{what} is
computed, not \emph{how}; each backend supplies its own bounds-safe implementation. Each
SYNAXIM hot op maps to exactly one opcode, so an entire transformer-replacement layer is a
short IR graph.

\section{Backend Lowering and Device Selection}\label{sec:backend}

AXIM lowers one IR to three families of target.

\textbf{CPU (SIMD).} The CPU backend emits vectorized kernels using AVX-512 or AVX2 on
x86-64 and NEON on ARM, with the instruction set chosen at \emph{runtime} by feature
detection. INT4 matvec unpacks nibbles with bitwise AND/SHIFT and accumulates in FP32.

\textbf{GPU (Vulkan / SPIR-V).} For NVIDIA, AMD, and Intel GPUs, compute kernels are authored
in GLSL, compiled to SPIR-V, and dispatched through Vulkan; the same SPIR-V module is
portable across all three vendors.

\textbf{GPU (Metal / MSL).} On Apple silicon, kernels are authored in MSL and compiled to a
\texttt{metallib}. This path is live and verified (\S\ref{sec:eval}).

\textbf{Device selection.} A program requests \texttt{device="auto"}, \texttt{"gpu"}, or
\texttt{"cpu"}. A Rust orchestrator (with a PyO3 bridge to the Python frontend) enumerates
devices and dispatches the lowered kernel. No path links or loads CUDA.

\textbf{CUDA porting shim (AXIM-HIP).} A header aliases common CUDA entry points (e.g.
\texttt{cudaMalloc}) onto AXIM calls, so a CUDA source can be recompiled against AXIM with a
single \texttt{\#include} and run on any AXIM device. This is a porting convenience, not a
runtime dependency.

\section{SYNAXIM as a Case Study}\label{sec:casestudy}

SYNAXIM is a standalone LLM inference engine that eliminates PyTorch, the Transformers
library, and the KV-Cache, replacing self-attention with a low-rank associative memory
$M = U V^\top$ and using an INT4 Boolean-bit pipeline for weight projections. Its
performance-critical operators are exactly the numerically sensitive, bandwidth-bound
kernels that must behave identically on every device. We express one decoder layer as AXIM
IR:

\begin{lstlisting}
RMSNORM -> INT4_MATVEC (Q/K/V, gate) -> LOWRANK_UPDATE -> LOWRANK_RETRIEVE
        -> RMSNORM -> SWIGLU (MLP) -> ADD (residual)
\end{lstlisting}

Each opcode lowers to the CPU-SIMD and GPU backends of \S\ref{sec:backend}. The critical
correctness requirement is that quantized projection produces the \emph{same} result on CPU
and GPU; \S\ref{sec:eval} shows this holds bit-for-bit on the verified platform.

\section{System Architecture}

\begin{table}[H]
\centering
\begin{tabular}{@{}lll@{}}
\toprule
Layer & Language & Role \\ \midrule
Frontend & Python & \texttt{@axim.kernel}, SYNAXIM ops, device API \\
IR & Python & typed buffers, operator graph \\
Orchestrator & Rust (+PyO3) & device enumeration and dispatch \\
CPU backend & C++ & AVX-512/AVX2/NEON + aximBLAS/aximDNN \\
GPU backend & C++/ObjC++ & Vulkan (SPIR-V) and Metal (MSL) dispatch \\
Graphics & ObjC++ (Metal) & render pipeline sharing the compute device \\
Shaders & GLSL / MSL & portable compute kernels \\ \bottomrule
\end{tabular}
\caption{AXIM layered architecture.}
\end{table}

Two tuned libraries accompany the CPU backend: \textbf{aximBLAS} (sgemv, sgemm, saxpy, sdot)
and \textbf{aximDNN} (softmax, layernorm, gelu, attention step)---CUDA-free analogues of
cuBLAS/cuDNN. A graphics pipeline renders on the same Metal device that runs compute,
enabling AI, physics, and rendering to share one queue. The implementation is ${\sim}3{,}400$
lines: 18 Python modules, 8 C++/Objective-C++ sources, 2 Rust sources, and 4 shader sources.

\section{Evaluation}\label{sec:eval}

We separate \textbf{verified} results (measured on hardware) from \textbf{projected} results.

\subsection{Verified on Apple~M3 (macOS, arm64)}

\begin{itemize}[leftmargin=1.4em]
    \item \textbf{CPU backend (NEON):} live; INT4 matvec, RMSNorm, SwiGLU, and low-rank
    retrieval execute correctly.
    \item \textbf{GPU backend (Metal):} live; the same operators execute on the M3 GPU.
    \item \textbf{Bit-exact cross-backend result:} for INT4 matvec, CPU (NEON) output equals
    GPU (Metal) output \emph{exactly}---the central correctness property for portable
    inference.
    \item \textbf{Full layer forward:} a complete SYNAXIM decoder layer runs end-to-end,
    CUDA-free.
    \item \textbf{Graphics:} the Metal render pipeline is live and shares the device with
    compute.
    \item \textbf{Automated tests:} 24 checks pass---pipeline/IR$+$dispatch (12), SYNAXIM
    ops (5), and the \texttt{.symb} loader (7).
\end{itemize}

\subsection{Projected (cross-vendor, Vulkan)}

The Vulkan path ships compiling SPIR-V compute shaders for the same operator set. Live
end-to-end execution on NVIDIA, AMD, and Intel GPUs requires the SPIR-V runtime executor
(loader), which is in progress. We therefore label cross-vendor GPU execution
\emph{projected}: the kernels exist and compile, but we do not yet report measured device
numbers. The CPU-SIMD path already compiles for x86-64 (AVX2) in continuous integration
alongside the macOS build.

\subsection{Discussion}

The verified bit-exact CPU\,$=$\,GPU result is the key evidence for AXIM's thesis: a single
IR, lowered independently to two unrelated backends (a scalar/SIMD CPU path and a Metal GPU
path), yields identical numerical output for quantized inference. This is precisely the
guarantee needed for ``write once, run anywhere'' model execution, and it holds without any
CUDA in the toolchain or runtime.

\section{Limitations and Future Work}

The Vulkan runtime executor (SPIR-V loader) that dispatches shaders on NVIDIA/AMD/Intel is in
progress; until then cross-vendor GPU results are projected. Publishing \texttt{maturin}
wheels that link the native backends is future work. A Windows (DirectX~12 / Vulkan) target
is planned. Zero-copy sharing of buffers between the render and compute queues (AI-in-games)
is designed but not yet wired end-to-end. Finally, aximBLAS/aximDNN cover the SYNAXIM path;
general HPC and graphics workloads will require an expanded kernel library.

\section{Conclusion}

AXIM shows that GPU and CPU compute need not be defined relative to CUDA. By compiling a
small, typed IR directly to vendor-neutral targets---SIMD on CPUs, Vulkan/Metal on GPUs---one
kernel definition runs across NVIDIA, AMD, Intel, and Apple silicon for AI, HPC, and graphics.
Using SYNAXIM as a case study, we verified on Apple~M3 that a complete post-transformer
inference layer runs CUDA-free with bit-exact CPU\,$=$\,GPU INT4 results. The runtime is
compact (${\sim}3{,}400$ lines), and its cross-vendor GPU path ships as compiling SPIR-V
today, with live execution as the immediate next milestone.

\begin{thebibliography}{10}
\bibitem{vaswani2017} Vaswani, A., et al. (2017). Attention Is All You Need. \emph{NeurIPS}.
\bibitem{vulkan} The Khronos Group. \emph{Vulkan} and \emph{SPIR-V} Specifications.
\bibitem{metal} Apple Inc. \emph{Metal Shading Language Specification}.
\bibitem{rocm} Advanced Micro Devices. \emph{ROCm and HIP Programming Guide}.
\bibitem{sycl} Intel Corporation. \emph{oneAPI / SYCL Specification}; SYCLomatic.
\bibitem{frantar2023} Frantar, E., et al. (2023). GPTQ. \emph{ICLR}.
\bibitem{lin2024} Lin, J., et al. (2024). AWQ. \emph{MLSys}.
\bibitem{xiao2023} Xiao, G., et al. (2023). SmoothQuant. \emph{ICML}.
\bibitem{synaxim2026} GRRNMAKER (2026). SYNAXIM: A Framework-Free LLM Inference Engine.
\url{https://github.com/GRRN-MAKER/SYNAXIM}.
\end{thebibliography}

\end{document}
